% ============================================================================
% APPA.TEX - Lampiran A: Instalasi Unity & Setup
% ============================================================================

\chapter{Instalasi Unity \& Setup}

\section{Download Unity Hub}

\begin{enumerate}
\item Kunjungi \url{https://unity.com/download}
\item Download Unity Hub untuk sistem operasi Anda
\item Install Unity Hub
\item Login atau buat akun Unity (gratis)
\end{enumerate}

\section{Install Unity Editor}

\begin{enumerate}
\item Buka Unity Hub
\item Klik "Installs" tab
\item Klik "Install Editor"
\item Pilih versi LTS (Long Term Support) terbaru
\item Pilih modules yang diperlukan:
   \begin{itemize}
   \item Windows Build Support (untuk Windows)
   \item Android Build Support (untuk Android)
   \item WebGL Build Support (untuk Web)
   \item Visual Studio Community (untuk coding)
   \end{itemize}
\item Klik "Install"
\end{enumerate}

\section{Membuat Project Baru}

\begin{enumerate}
\item Buka Unity Hub
\item Klik "New Project"
\item Pilih template:
   \begin{itemize}
   \item 3D: Untuk game 3D
   \item 2D: Untuk game 2D
   \item VR: Untuk VR projects
   \end{itemize}
\item Tentukan nama project
\item Pilih lokasi folder
\item Klik "Create Project"
\end{enumerate}

\section{Konfigurasi Awal}

\subsection{Project Settings}

\begin{itemize}
\item Edit > Project Settings > Player
\item Set Company Name dan Product Name
\item Configure Icon dan Splash Screen
\item Set Default Resolution
\end{itemize}

\subsection{Input Settings}

\begin{itemize}
\item Edit > Project Settings > Input Manager
\item Configure input axes
\item Set up keyboard, mouse, gamepad inputs
\end{itemize}

\section{Troubleshooting}

\subsection{Unity Tidak Bisa Dibuka}

\begin{itemize}
\item Pastikan sistem memenuhi requirements
\item Update graphics drivers
\item Check antivirus tidak memblokir Unity
\end{itemize}

\subsection{Project Tidak Bisa Dibuka}

\begin{itemize}
\item Pastikan versi Unity sesuai
\item Check file project tidak corrupt
\item Reimport project jika perlu
\end{itemize}

\section{Resources Tambahan}

\begin{itemize}
\item Unity Learn: \url{https://learn.unity.com/}
\item Unity Documentation: \url{https://docs.unity3d.com/}
\item Unity Forum: \url{https://forum.unity.com/}
\end{itemize}
